\begin{textsample}{POS Dim 4 – rs – Score 6.00 – rs/rs\_em\_12.txt}  \label{ex:f4_pos_105}
1.7.4 Educação Escolar \textbf{Indígena} e Educação Escolar Quilombola As modalidades de educação escolar \textbf{indígena} e educação escolar quilombola são tratadas a partir do respeito, do reconhecimento e do valor ético-cultural e como constituintes do sistema estadual e municipais de educação do Rio Grande do Sul.
Este RCGEM apresenta legislações específicas para subsidiar a implementação e a consolidação dessas modalidades, na etapa final da Educação Básica, e desenvolver as propostas pedagógicas demandadas pelas \textbf{populações} \textbf{indígenas} e quilombolas consideradas fundamentais para a compreensão dos direitos educacionais específicos.
É imprescindível que as redes de ensino que compõem o \textbf{território} gaúcho conheçam a legislação, a estudem e a apliquem no trabalho pedagógico-educativo.
O poder público, em suas instâncias administrativas, observando o estabelecido na Constituição Federal e Estadual, assume o compromisso de viabilizar, financiar, apoiar e desenvolver as modalidades escolares de educação \textbf{indígena} e quilombola.
As propostas pedagógicas devem contemplar a legislação internacional, nacional e gaúcha nas suas constituições, arquiteturas, metodologias, manutenção, desenvolvimento e financiamento.
Para ressaltar essa necessidade, o RCGEM apresenta aspectos da legislação que devem embasar as propostas pedagógicas das modalidades de educação escolar \textbf{indígena} e quilombola, sem prejuízos das autonomias dos povos específicos, compreendidos como integrantes das sociedades brasileira e gaúcha.
Portanto, são legislações pertinentes, centrais, fundamentais, sobre as quais devem ser realizados estudos consistentes.
A Convenção nº 169/1989 da Organização Internacional do Trabalho ( OIT, 1989 ), relativa aos direitos dos povos \textbf{indígenas} e das comunidades quilombolas, prevê programas e serviços de educação a serem desenvolvidos e aplicados em cooperação com os povos específicos, para responder às suas necessidades particulares.
Essas ações devem abranger a sua história, seus conhecimentos e técnicas, seus sistemas de valores e todas suas aspirações sociais, econômicas e culturais.
O documento estabelece também o direito e a garantia de formação de membros dos povos \textbf{indígenas} e quilombolas na formulação e execução de programas de educação, com vistas a lhes transferir, progressivamente, as responsabilidades de suas realizações.
A mesma Convenção Internacional compromete os governos ao reconhecimento dos direitos dos \textbf{indígenas} e quilombolas a criarem suas próprias instituições e meios de educação em consonância com as normas definidas coletivamente entre os povos e a autoridade competente.
O RCGEM, amparado pela legislação e normatização internacional, nacional e do \textbf{território} gaúcho, estabelece a necessária observância, na elaboração de documentos e normas educacionais que envolvam as línguas, os saberes e as pedagogias dos povos \textbf{indígenas} e das comunidades quilombolas, no núcleo comum, seja realizado pelo povo \textbf{indígena} ( Kaingang, Guarani, Xokleng e Charrua ) e pelas comunidades quilombolas.
Desse modo, como prevê a legislação e este RCGEM, os povos e comunidades \textbf{indígenas} e quilombolas possam eles próprios detalhar em suas línguas ( com tradução ) os princípios educativos de suas culturas e pedagogias.
Implanta -se, com essa metodologia, um diálogo intercultural equânime na Educação Básica com a BNCC e a comunidade brasileira e gaúcha, autenticidade, ética e propriedade na composição dos projetos educacionais e curriculares, especificamente, \textbf{indígenas} e quilombolas.

% matched lemmas: indígena, população, território
\end{textsample}
